\Needspace{12\baselineskip}
\begin{center}\emph{Strenger Beweis des Fundamentalsatzes.}\end{center}

\Article{135}
Im vorigen Artikel wurden die Aussagen aus Artikel $133$ zwar nicht bewiesen. Es wurde aber gezeigt, dass ihre Richtigkeit von der des Fundamentalsatzes abhängt, die wir vorläufig vorausgesetzt haben. Aus der Herleitung selbst geht jedoch hervor, dass diese Aussagen für die Zahlen $P,Q$ gelten, sobald der Fundamentalsatz beim Vergleich aller Primfaktoren dieser Zahlen miteinander gilt, selbst wenn er nicht allgemein richtig ist. Wir gehen nun zum Beweis des Fundamentalsatzes selbst über und stellen ihm die folgende Erklärung voran.

\emph{Wir sagen, der Fundamentalsatz sei bis zu einer Zahl $M$ richtig, wenn er für zwei beliebige Primzahlen gilt, von denen keine $M$ übersteigt.}

Entsprechend ist es zu verstehen, wenn wir sagen, dass die Sätze aus den Artikeln $131,132,133$ bis zu einer bestimmten Grenze richtig sind. Ist der Fundamentalsatz bis zu einer Grenze als richtig bekannt, so gelten offensichtlich auch diese Aussagen bis zu derselben Grenze.

\SourceLeaf{110}\Article{136}
Für kleine Zahlen lässt sich der Fundamentalsatz leicht durch Nachprüfen bestätigen. So erhält man eine Grenze, bis zu der er sicher gilt. Wir setzen voraus, dass diese Prüfung durchgeführt wurde; wie weit wir sie fortgesetzt haben, spielt aber keine Rolle. Es würde sogar genügen, wenn wir ihn nur bis zur Zahl $5$ bestätigt hätten. Dazu reicht die einzige Beobachtung, dass $+5N3$, $\pm3N5$ gilt.

Ist nun der Fundamentalsatz nicht allgemein richtig, so gibt es eine Grenze $T$, bis zu der er gilt, während er bis zur nächstgrößeren Zahl $T+1$ nicht mehr gilt. Das heißt: Es gibt zwei Primzahlen, deren größere $T+1$ ist und deren Vergleich dem Fundamentalsatz widerspricht; je zwei beliebige andere Primzahlen stimmen dagegen mit ihm überein, sofern beide kleiner als $T+1$ sind. Also gelten auch die Aussagen aus den Artikeln $131,132,133$ bis zu $T$. Wir zeigen nun, dass diese Annahme unmöglich ist. Dabei müssen wir die folgenden Fälle unterscheiden, je nachdem, welche Formen $T+1$ und die kleinere Primzahl haben, die kleiner als $T+1$ ist und deren Vergleich mit $T+1$ nach Voraussetzung dem Satz widerspricht. Wir bezeichnen diese Primzahl mit $p$.

Haben sowohl $T+1$ als auch $p$ die Form $4n+1$, so könnte der Fundamentalsatz auf zwei Arten falsch sein, nämlich wenn gleichzeitig \emph{entweder}

\[
\begin{array}{ll}&\pm pR(T+1)\quad\text{und}\quad\pm(T+1)Np\\\textit{oder zugleich}&\pm pN(T+1)\quad\text{und}\quad\pm(T+1)Rp\end{array}
\]

Haben sowohl $T+1$ als auch $p$ die Form $4n+3$, so wäre der Fundamentalsatz falsch, wenn gleichzeitig \emph{entweder}

\[
\begin{array}{ll}&+pR(T+1)\quad\text{und}\quad-(T+1)Np\\\text{(gleichbedeutend}&-pN(T+1)\quad\text{und}\quad+(T+1)Rp)\\\textit{oder}&+pN(T+1)\quad\text{und}\quad-(T+1)Rp\\\text{(oder}&-pR(T+1)\quad\text{und}\quad+(T+1)Np)\end{array}
\]

Hat $T+1$ die Form $4n+1$, $p$ dagegen die Form $4n+3$, so wäre der Fundamentalsatz falsch, wenn \emph{entweder}

\[
\begin{array}{ll}&\pm pR(T+1)\ \text{und}\ +(T+1)Np\ (\text{oder}\ -(T+1)Rp)\\\textit{oder}&\pm pN(T+1)\ \text{und}\ -(T+1)Np\ (\text{oder}\ +(T+1)Rp)\end{array}
\]

\SourceLeaf{111}Hat $T+1$ die Form $4n+3$, $p$ dagegen die Form $4n+1$, so wäre der Fundamentalsatz falsch, wenn \emph{entweder}

\[
\begin{array}{ll}&+pR(T+1)\ (\text{oder}\ -pN(T+1))\ \text{und}\ \pm(T+1)Np\\\textit{oder}&+pN(T+1)\ (\text{oder}\ -pR(T+1))\ \text{und}\ \pm(T+1)Rp\end{array}
\]

Können wir beweisen, dass keiner dieser acht Fälle eintritt, so steht zugleich fest, dass die Richtigkeit des Fundamentalsatzes durch keine Grenze beschränkt ist. Diese Aufgabe gehen wir nun an. Da aber einige dieser Fälle von anderen abhängen, können wir die Reihenfolge, in der wir sie hier aufgezählt haben, nicht beibehalten.

\Needspace{9\baselineskip}
\Article{137}
\emph{Erster Fall. Hat $T+1$ die Form $4n+1$ $(=a)$, und $p$ dieselbe Form, und gilt außerdem $\pm pRa$, so kann nicht $\pm aNp$ gelten.} Dies war oben der erste Fall.

Sei $+p\equiv e^2(\mathrm{mod.}\ a)$, wobei $e$ gerade und $<a$ ist (was sich immer erreichen lässt). Nun sind zwei Fälle zu unterscheiden.

I.\quad Ist $e$ nicht durch $p$ teilbar, so setze man $e^2=p+af$. Dann ist $f$ positiv, von der Form $4n+3$ (oder der Form $B$), $<a$ und nicht durch $p$ teilbar. Ferner gilt $e^2\equiv p(\mathrm{mod.}\ f)$, \emph{das heißt} $pRf$, also nach Aussage $11$ aus Artikel $132$ auch $\pm fRp$ (denn wegen $p,f<a$ gelten diese Aussagen für sie). Es gilt aber auch $afRp$, also ebenfalls $\pm aRp$.

II.\quad Ist $e$ durch $p$ teilbar, so setze man $e=gp$ und $e^2=p+aph$ oder $pg^2=1+ah$. Dann hat $h$ die Form $4n+3$ ($B$) und ist zu $p$ und $g^2$ teilerfremd. Ferner gilt $pg^2Rh$, also auch $pRh$; nach Aussage $11$ aus Artikel $132$ folgt daraus $\pm hRp$. Es gilt aber auch $-ahRp$, da $-ah\equiv1(\mathrm{mod.}\ p)$, also ebenfalls $\mp aRp$.

\Needspace{9\baselineskip}
\Article{138}
\emph{Zweiter Fall. Hat $T+1$ die Form $4n+1$ $(=a)$, $p$ die Form $4n+3$, und gilt $\pm pR(T+1)$, so kann nicht $+(T+1)Np$ oder $-(T+1)Rp$ gelten.} Dies war oben der fünfte Fall.

Wie oben sei $e^2=p+fa$, wobei $e$ gerade und $<a$ ist.

I.\quad Ist $e$ nicht durch $p$ teilbar, so ist auch $f$ nicht durch $p$ teilbar. Außerdem ist $f$ positiv, von der Form $4n+1$ (oder $A$) und $<a$. Es gilt $+pRf$, \SourceLeaf{112}also nach Aussage 10 aus Artikel 132 auch $+fRp$. Es gilt aber auch $+faRp$, also $+aRp$ oder $-aNp$.

II. Ist $e$ durch $p$ teilbar, so sei $e=pg$ und $f=ph$. Somit gilt $g^2p=1+ha$. Dann ist $h$ positiv, von der Form $4n+3$ $(B)$ und zu $p$ und $g^2$ teilerfremd. Ferner gilt $+g^2pRh$, also $+pRh$; daraus folgt nach Aussage 13 aus Artikel 132 $-hRp$. Es gilt aber $-haRp$, woraus $+aRp$ und $-aNp$ folgt.

\Needspace{10\baselineskip}
\Article{139}
\textit{Dritter Fall. Hat $T+1$ die Form $4n+1$ $(=a)$, $p$ dieselbe Form, und gilt $\pm pNa$, so kann nicht $\pm aRp$ gelten.} (Oben der zweite Fall.)

Man nehme eine Primzahl unterhalb von $a$, von der $+a$ ein Nichtrest ist. Dass es solche Primzahlen gibt, haben wir oben bewiesen (Artikel 125, 129). Hier sind aber zwei Fälle gesondert zu betrachten, je nachdem, ob diese Primzahl die Form $4n+1$ oder $4n+3$ hat. Denn es wurde nicht bewiesen, dass es solche Primzahlen \textit{beider} Formen gibt.

I. Diese Primzahl habe die Form $4n+1$ und sei $=a'$. Dann gilt $\pm a'Na$ (Artikel 131), also $\pm a'pRa$. Sei daher $e^2\equiv a'p\ (\mathrm{mod.}\ a)$, wobei $e$ gerade und $<a$ ist. Nun sind wieder vier Fälle zu unterscheiden.

1) Ist $e$ weder durch $p$ noch durch $a'$ teilbar, so setze man $e^2=a'p\pm af$ und wähle die Vorzeichen so, dass $f$ positiv ist. Dann gilt $f<a$, und die Zahl ist zu $a'$ und $p$ teilerfremd; beim oberen Zeichen hat sie die Form $4n+3$, beim unteren die Form $4n+1$. Der Kürze wegen bezeichnen wir mit $[x,y]$ die Anzahl der Primfaktoren der Zahl $y$, von denen $x$ ein Nichtrest ist. Dann gilt $a'pRf$, also $[a'p,f]=0$. Somit ist $[f,a'p]$ gerade (Aussagen 1, 3 aus Artikel 133), \textit{das heißt} entweder $=0$ oder $=2$. Daher ist $f$ entweder ein Rest beider Zahlen $a',p$ oder keiner von beiden. Ersteres ist unmöglich, denn $\pm af$ ist ein Rest von $a'$, und nach Voraussetzung gilt $\pm aNa'$; daraus folgt $\pm fNa'$. Also muss $f$ ein Nichtrest beider Zahlen $a',p$ sein. Wegen $\pm afRp$ gilt dann $\pm aNp$. \mbox{\textit{Q. E. D.}}

2) Ist $e$ durch $p$, aber nicht durch $a'$ teilbar, so sei $e=gp$ und $g^2p=a'\pm ah$, wobei das Zeichen so gewählt wird, dass $h$ positiv ist. Dann gilt $h<a$, und die Zahl ist zu $a',g$ und $p$ teilerfremd; beim oberen Zeichen hat sie die Form $4n+3$, beim unteren die Form $4n+1$. Aus der Gleichung $g^2p=a'\pm ah$ erhält man durch Multiplikation mit $p$ und $a'$ ohne Schwierigkeit $pa'Rh\,.\,.\,.\,.\,.(\alpha)$; $\pm ahpRa'\,.\,.\,.\,.(\beta)$; $aa'hRp\,.\,.\,.\,.\,.(\gamma)$. Aus $(\alpha)$ folgt $[pa',h]=0$, also ist (nach den Aussagen 1, 3 aus Artikel 133) $[h,pa']$ gerade, \textit{das heißt,} \SourceLeaf{113}$h$ ist entweder ein Nichtrest beider Zahlen $p,a'$ oder keiner von beiden. \textit{Im ersten Fall} folgt aus $(\beta)$ $\pm apNa'$. Da nach Voraussetzung $\pm aNa'$ gilt, folgt $\pm pRa'$. Nach dem Fundamentalsatz, der für die Zahlen $p,a'$ unterhalb von $T+1$ gilt, ergibt sich daraus $\pm a'Rp$. Hieraus und aus $hNp$ folgt nach $(\gamma)$ $\pm aNp$. \mbox{\textit{Q. E. D.}} \textit{Im zweiten Fall} folgt aus $(\beta)$ $\pm apRa'$, daraus $\pm pNa'$, $\pm a'Np$ und schließlich zusammen mit $hRp$ nach $(\gamma)$ $\pm aNp$. \mbox{\textit{Q. E. D.}}

3) Ist $e$ durch $a'$, aber nicht durch $p$ teilbar, so verläuft der Beweis fast genauso wie im vorigen Fall. Er wird niemanden aufhalten, der jenen gründlich verstanden hat.

{\widowpenalty=10000 4) Ist $e$ sowohl durch $a'$ als auch durch $p$ und damit auch durch ihr Produkt $a'p$ teilbar (wir nehmen die Zahlen $a',p$ nämlich als \textit{verschieden} an, da sonst das zu Beweisende, nämlich $aNp$, schon in der Voraussetzung $aNa'$ enthalten wäre), so sei $e=ga'p$ und $g^2a'p=1\pm ah$. Dann gilt $h<a$, und die Zahl ist zu $a'$ und $p$ teilerfremd; beim oberen Zeichen hat sie die Form $4n+3$, beim unteren die Form $4n+1$. Aus dieser Gleichung lassen sich offensichtlich die folgenden Beziehungen herleiten: $a'pRh\,.\,.\,.\,.(\alpha)$; $\pm ahRa'\,.\,.\,.\,.(\beta)$; $\pm ahRp\,.\,.\,.\,.(\gamma)$. Aus $(\alpha)$, das mit $(\alpha)$ in (2) übereinstimmt, folgt wie dort, dass entweder gleichzeitig $hRp$, $hRa'$ oder $hNp$, $hNa'$ gilt. Im ersten Fall würde aber nach $(\beta)$ $aRa'$ gelten, entgegen der Voraussetzung. Also gilt $hNp$, nach $(\gamma)$ somit auch $aNp$.\par}

II. Hat diese Primzahl die Form $4n+3$, so ist der Beweis dem vorigen so ähnlich, dass wir es für überflüssig hielten, ihn anzugeben. Für diejenigen, die ihn selbst ausarbeiten möchten (was wir sehr empfehlen), bemerken wir nur: Ist man zu einer Gleichung wie $e^2=bp\pm af$ gelangt (wobei $b$ jene Primzahl bezeichnet), so wird der Beweis übersichtlicher, wenn man beide Vorzeichen gesondert betrachtet.

\Needspace{8\baselineskip}
\Article{140}
\textit{Vierter Fall. Hat $T+1$ die Form $4n+1$ $(=a)$, $p$ die Form $4n+3$, und gilt $\pm pNa$, so kann nicht $+aRp$ oder $-aNp$ gelten.} (Oben der sechste Fall.)

Auch den Beweis dieses Falles lassen wir der Kürze wegen weg, da er genauso verläuft wie der Beweis des dritten Falles.

\SourceLeaf{114}\Needspace{9\baselineskip}
\Article{141}
\textit{Fünfter Fall. Hat $T+1$ die Form $4n+3$ $(=b)$, $p$ dieselbe Form, und gilt $+pRb$ oder $-pNb$, so kann nicht $+bRp$ oder $-bNp$ gelten.} (Oben der dritte Fall.)

Sei $p\equiv e^2\ (\mathrm{mod.}\ b)$, wobei $e$ gerade und $<b$ ist.

I. Ist $e$ nicht durch $p$ teilbar, so setze man $e^2=p+bf$. Dann ist $f$ positiv, von der Form $4n+3$, $<b$ und zu $p$ teilerfremd. Ferner gilt $pRf$, also nach Aussage 13 aus Artikel 132 $-fRp$. Hieraus und aus $+bfRp$ folgt $-bRp$, also $+bNp$. \mbox{\textit{Q. E. D.}}

II. Ist $e$ durch $p$ teilbar, so sei $e=pg$ und $ggp=1+bh$. Dann hat $h$ die Form $4n+1$ und ist zu $p$ teilerfremd; es gilt $p\equiv g^2p^2\ (\mathrm{mod.}\ h)$, also $pRh$. Daraus folgt $+hRp$ (Aussage 10 aus Artikel 132), und hieraus zusammen mit $-bhRp$ ergibt sich $-bRp$ oder $+bNp$. \mbox{\textit{Q. E. D.}}

\Needspace{8\baselineskip}
\Article{142}
\textit{Sechster Fall. Hat $T+1$ die Form $4n+3$ $(=b)$, $p$ die Form $4n+1$, und gilt $pRb$, so kann nicht $\pm bNp$ gelten.} (Oben der siebente Fall.)

Den Beweis, der genauso wie der vorige verläuft, lassen wir weg.

\Needspace{9\baselineskip}
\Article{143}
\textit{Siebenter Fall. Hat $T+1$ die Form $4n+3$ $(=b)$, $p$ dieselbe Form, und gilt $+pNb$ oder $-pRb$, so kann nicht $+bNp$ oder $-bRp$ gelten.} (Oben der vierte Fall.)

Sei $-p\equiv e^2\ (\mathrm{mod.}\ b)$, wobei $e$ gerade und $<b$ ist.

I. Ist $e$ nicht durch $p$ teilbar, so sei $-p=e^2-bf$. Dann ist $f$ positiv, von der Form $4n+1$, zu $p$ teilerfremd und kleiner als $b$ (denn $e$ ist sicher nicht größer als $b-1$, und $p<b-1$, also gilt $bf=e^2+p<b^2-b$, \textit{das heißt} $f<b-1$). Ferner gilt $-pRf$, nach Aussage 10 aus Artikel 132 somit $+fRp$; hieraus und aus $+bfRp$ folgt $+bRp$ oder $-bNp$.

II. Ist $e$ durch $p$ teilbar, so sei $e=pg$ und $g^2p=-1+bh$. Dann ist $h$ positiv, von der Form $4n+3$, zu $p$ teilerfremd und $<b$. Ferner gilt $-pRh$, nach Aussage 14 aus Artikel 132 also $+hRp$. Hieraus und aus $bhRp$ folgt $+bRp$ oder $-bNp$. \mbox{\textit{Q. E. D.}}

\SourceLeaf{115}\Needspace{8\baselineskip}
\Article{144}
\textit{Achter Fall. Hat $T+1$ die Form $4n+3$ $(=b)$, $p$ die Form $4n+1$, und gilt $+pNb$ oder $-pRb$, so kann nicht $\pm bRp$ gelten.} (Oben der letzte Fall.)

Der Beweis verläuft wie im vorigen Fall.
